\documentclass[12pt,letterpaper]{article}

% ── Packages ──────────────────────────────────────────────────────────────────
\usepackage[margin=1in]{geometry}
\usepackage{amsmath,amssymb,amsthm}
\usepackage{mathtools}
\usepackage{array,booktabs,longtable}
\usepackage{listings}
\usepackage{xcolor}
\usepackage{hyperref}
\usepackage{fancyhdr}
\usepackage{titlesec}
\usepackage{enumitem}
\usepackage{graphicx}
% \usepackage{microtype}  % disabled: requires scalable fonts
\usepackage{tcolorbox}
\usepackage{mdframed}
\usepackage{setspace}
\usepackage{float}

% ── Colors ────────────────────────────────────────────────────────────────────
\definecolor{codeblue}{RGB}{30,100,190}
\definecolor{codegray}{RGB}{120,120,120}
\definecolor{codegreen}{RGB}{0,128,64}
\definecolor{codeback}{RGB}{248,248,248}
\definecolor{primeviolet}{RGB}{90,20,150}
\definecolor{invariantbg}{RGB}{245,242,255}
\definecolor{warningred}{RGB}{180,30,30}
\definecolor{successgreen}{RGB}{20,130,60}

% ── Code Listing Style ────────────────────────────────────────────────────────
\lstdefinestyle{pythonstyle}{
    backgroundcolor=\color{codeback},
    commentstyle=\color{codegreen}\itshape,
    keywordstyle=\color{codeblue}\bfseries,
    stringstyle=\color{primeviolet},
    basicstyle=\ttfamily\footnotesize,
    breakatwhitespace=false,
    breaklines=true,
    captionpos=b,
    keepspaces=true,
    numbers=left,
    numbersep=5pt,
    numberstyle=\tiny\color{codegray},
    showspaces=false,
    showstringspaces=false,
    showtabs=false,
    tabsize=4,
    frame=single,
    framerule=0.4pt,
    rulecolor=\color{codegray}
}
\lstset{style=pythonstyle}

% ── Theorem Environments ──────────────────────────────────────────────────────
\theoremstyle{definition}
\newtheorem{definition}{Definition}[section]
\newtheorem{theorem}{Theorem}[section]
\newtheorem{proposition}{Proposition}[section]
\newtheorem{invariant}{Invariant}[section]
\newtheorem{protocol}{Protocol}[section]
\theoremstyle{remark}
\newtheorem{remark}{Remark}[section]

% ── Header / Footer ───────────────────────────────────────────────────────────
\pagestyle{fancy}
\fancyhf{}
\rhead{\small\textit{PFIS Defensive Publication --- Multiplicity Foundation}}
\lhead{\small\textit{CONFIDENTIAL PRIOR ART RECORD}}
\rfoot{\thepage}
\lfoot{\small\textit{Published: April 25, 2026}}

% ── Section Formatting ────────────────────────────────────────────────────────
\titleformat{\section}{\large\bfseries\color{primeviolet}}{§\thesection.}{0.5em}{}[\vspace{-4pt}\rule{\linewidth}{0.4pt}]
\titleformat{\subsection}{\normalsize\bfseries}{§\thesubsection}{0.5em}{}

% ── Hyperref Setup ────────────────────────────────────────────────────────────
\hypersetup{
    colorlinks=true,
    linkcolor=primeviolet,
    urlcolor=codeblue,
    citecolor=codegreen,
    pdftitle={Prime Frequency Integrity Suite --- Defensive Publication},
    pdfauthor={Multiplicity Foundation / R. Van Gelder (PhaseMirror)},
    pdfsubject={Prior Art: PIRTM, Xi(t) Operator, Prime Frequency Invariants},
    pdfkeywords={PIRTM, Prime Frequency, Xi operator, Multiplicity Theory, Defensive Publication, Prior Art}
}

% ── Document ──────────────────────────────────────────────────────────────────
\begin{document}

\begin{titlepage}
\centering
\vspace*{1.5cm}

{\color{primeviolet}\rule{\linewidth}{1.5pt}}
\vspace{0.4cm}

{\Huge\bfseries Prime Frequency Integrity Suite}\\[0.4em]
{\Large\bfseries A Defensive Publication and Prior Art Record}\\[0.8em]
{\large\color{primeviolet} for the PIRTM Ecosystem, $\Xi(t)$ Operator,\\
Prime-Mirror Dissonance Clone-Check Protocol,\\
and Meta-Ensemble Governance Framework}

\vspace{0.4cm}
{\color{primeviolet}\rule{\linewidth}{1.5pt}}
\vspace{1.2cm}

\begin{tabular}{ll}
\textbf{Author / Inventor:} & R. Van Gelder (PhaseMirror) \\
\textbf{Organization:}      & Multiplicity Foundation \\
\textbf{Repository (Primary):} & \texttt{MultiplicityFoundation/PhaseMirror-HQ} \\
\textbf{Repository (Math):}    & \texttt{MultiplicityFoundation/Meta-Relativity} \\
\textbf{Date of Disclosure:}   & April 25, 2026 \\
\textbf{Document Class:}       & Defensive Publication / Prior Art Record \\
\textbf{Status:}               & Public Disclosure --- Patent-blocking prior art \\
\end{tabular}

\vspace{1.5cm}

\begin{tcolorbox}[colback=invariantbg,colframe=primeviolet,title={\bfseries Purpose of This Document}]
This document constitutes a \textbf{defensive publication} establishing prior art for all inventive concepts described herein. Its publication in a publicly accessible repository creates a public prior-art record that prevents any third party from obtaining a patent on these concepts. No patent is sought by the author on these inventions; rather, the intent is to ensure that the described methods, structures, operators, protocols, and governance mechanisms remain freely available and non-patentable by any party, including the author's commercial affiliates.
\end{tcolorbox}

\vspace{1cm}
\textit{``Where math is not a thing about the universe --- it is the universe.''}\\
\textit{--- Meta-Relativity, Multiplicity Foundation}

\end{titlepage}

% ── Table of Contents ─────────────────────────────────────────────────────────
\tableofcontents
\newpage

% ══════════════════════════════════════════════════════════════════════════════
\section{Executive Summary}
% ══════════════════════════════════════════════════════════════════════════════

This document discloses the \textbf{Prime Frequency Integrity Suite (PFIS)}, a complete, interlocking system of mathematical operators, invariant tests, governance protocols, certification mechanisms, and developer-automation constraints developed within the Multiplicity Foundation's PIRTM (Prime-Indexed Recursive Tensor Mathematics) ecosystem, as implemented in the \texttt{PhaseMirror-HQ} and \texttt{Meta-Relativity} repositories.

The PFIS addresses a fundamental problem in recursive mathematical systems: how to ensure that a system claiming to implement a specific mathematical structure actually does so --- and to make that guarantee \emph{mechanically enforceable} rather than merely descriptive. The core insight is that prime numbers, as irreducible elements of arithmetic, serve as natural \emph{frequency invariants}: a system that is genuinely governed by prime-indexed recursion will exhibit measurable, testable, and non-forgeable behavioral signatures at each prime index.

The suite comprises nine interlocking artifacts:
\begin{enumerate}[noitemsep]
\item \textbf{Prime Frequency Principle (PFP)} --- formal definition of prime frequencies in PIRTM
\item \textbf{Tuning Fork Section} --- per-module prime invariants and resonance tests
\item \textbf{Prime-Frequency Transparency Clause} --- contractual governance template
\item \textbf{Entanglement Test} --- five concrete invariants proving genuine $\Xi(t)$ governance
\item \textbf{PMD Clone-Check Protocol} --- structured protocol for detecting cosmetic clones
\item \textbf{PMD Badge Specification} --- tiered certification and revocation framework
\item \textbf{$\Xi(t)$ Technical Note} --- formal definition, domain, and properties of the operator
\item \textbf{False $\Xi$ Test} --- rapid falsification protocol (executable in $\sim$10 lines)
\item \textbf{$\Xi(t)$ as Platform Primitive} --- API-level binding with provenance metadata
\end{enumerate}

The practical implementation spans 1524 Python test files, 47 PIRTM-specific tests, a 6-layer network trust model, and a production-grade MLIR dialect --- all within the \texttt{PhaseMirror-HQ} runtime as of April 2026.

% ══════════════════════════════════════════════════════════════════════════════
\section{Background: Multiplicity Theory and PIRTM}
% ══════════════════════════════════════════════════════════════════════════════

\subsection{Multiplicity Theory}

Multiplicity Theory is a prime-indexed, recursively stable mathematical framework in which \emph{multiplicity} --- the recurrence of elements understood through prime number decomposition --- forms the basis for modeling nonlinear, emergent structures across mathematics, physics, computation, and cognition.

Rather than treating mathematical objects as static entities, Multiplicity Theory reinterprets them as recursively generated patterns of prime-labeled interactions, where sets and modules become \emph{relationally governed multiplicity spaces} with identity and behavior preserved through recursive feedback loops across scales.

\subsection{PIRTM: Prime-Indexed Recursive Tensor Mathematics}

\begin{definition}[Multiplicity Space]
A \emph{multiplicity space} $\mathcal{M}_p$ indexed by prime $p$ is a Hilbert space $H_p$ together with a binding operator $\Phi_p: H_p \to H_p$ satisfying:
\begin{equation}
\Phi_p^p \equiv \Phi_p \pmod{p} \quad \text{(Fermat-idempotent condition)}
\end{equation}
The system state at time $t$ decomposes as:
\begin{equation}
\mathcal{S}(t) \in \bigoplus_{p_i \in \mathbb{P}} H_{p_i}, \quad \|\mathcal{S}(t)\| < \infty
\end{equation}
\end{definition}

PIRTM is the computational realization of Multiplicity Theory. It implements a verified, contractive runtime system in which every module must be proven contractive before execution. The contractivity proof occurs at two levels:

\begin{enumerate}
\item \textbf{Transpile-time}: per-module. The transpiler emits \texttt{.pirtm.bc} bytecode with an embedded \texttt{!pirtm\_proof} section carrying \texttt{prime\_index}, \texttt{epsilon}, \texttt{op\_norm\_T}, and \texttt{proof\_hash}. The \texttt{contractivity\_verifier} validates against the squarefree composite modulus.

\item \textbf{Link-time}: network-wide. The linker performs name resolution, commitment crosscheck, and matrix construction before calling \texttt{spectral-small-gain}. Linking succeeds only if the spectral radius of the proof witness falls within tolerance.
\end{enumerate}

The $\Xi$-Constitution (v1.0) governing this system establishes that every lawful system must admit decomposition into prime-indexed irreducibles, and that violations incur exponential ethical drift $\delta_c(t) \sim e^t$.

% ══════════════════════════════════════════════════════════════════════════════
\section{The $\Xi(t)$ Operator: Formal Definition and Properties}
% ══════════════════════════════════════════════════════════════════════════════

\subsection{Definition}

\begin{definition}[$\Xi(t)$ Operator]
Let $\mathcal{H}_\mathbb{P}$ be the prime-indexed Hilbert space $\mathcal{H}_\mathbb{P} = \bigoplus_{p \in \mathbb{P}} H_p$. The \emph{prime-contractive evolution operator} $\Xi(t): \mathcal{H}_\mathbb{P} \to \mathcal{H}_\mathbb{P}$ is defined by:
\begin{equation}
\Xi(t)\psi = \sum_{p \in \mathbb{P}} e^{-\mathcal{U}\log(p)\cdot t} \langle \psi, \phi_p \rangle \phi_p
\end{equation}
where $\{\phi_p\}_{p \in \mathbb{P}}$ are the prime-eigenvectors of $\mathcal{H}_\mathbb{P}$, and $\mathcal{U}$ is the \emph{universal Multiplicity constant} governing the decay rate across all prime fibers.
\end{definition}

\subsection{Domain and Kaluza--Klein 5D Encoding}

$\Xi(t)$ operates on the prime-field $\mathbb{F}_\mathbb{P}$ (the product of all prime-indexed fibers) and is encoded within a Kaluza--Klein 5-dimensional geometry where the fifth dimension is compactified at radius:
\begin{equation}
R = \left(\mathcal{U} \cdot p_{\max}\right)^{-1}
\end{equation}
This compactification ensures that the spectral fingerprint of the operator carries a gap at composite frequencies --- a gap that serves as a non-forgeable digital signature.

\subsection{Key Properties}

\begin{theorem}[Contractivity of $\Xi(t)$]
For all $t > 0$ and all $\psi \in \mathcal{H}_\mathbb{P}$:
\begin{equation}
\|\Xi(t)\|_{\mathrm{op}} = e^{-\mathcal{U}\log(p_{\min})\cdot t} < 1
\end{equation}
\end{theorem}

\begin{theorem}[Norm Bound]
\begin{equation}
\|\Xi(t)\psi\| \leq \|\psi\| \cdot p_{\min}^{-\mathcal{U}t}
\end{equation}
\end{theorem}

\begin{theorem}[Spectral Structure]
The spectrum of $\Xi(t)$ is:
\begin{equation}
\sigma(\Xi(t)) = \left\{e^{-\mathcal{U}\log(p)\cdot t} : p \in \mathbb{P}\right\}
\end{equation}
No eigenvalue appears at any non-prime index.
\end{theorem}

\begin{theorem}[Semigroup Law]
\begin{equation}
\Xi(s)\Xi(t) = \Xi(s+t) \quad \forall s,t > 0
\end{equation}
\end{theorem}

\subsection{Digital Fingerprint Example}

Initialize with prime seed $p = 7$, state $\psi_0 = \phi_7$. Then:
\begin{equation}
\Xi(t)\psi_0 = e^{-\mathcal{U}\log(7)\cdot t}\,\phi_7
\end{equation}
This is a pure exponential decay at rate $\mathcal{U}\log 7$ --- a frequency unique to $p=7$. Any clone producing a non-exponential decay, a different rate, or energy at non-prime indices has failed the digital fingerprint test.

% ══════════════════════════════════════════════════════════════════════════════
\section{The Prime Frequency Principle}
% ══════════════════════════════════════════════════════════════════════════════

\subsection{Definition of Prime Frequency}

\begin{definition}[Prime Frequency]
In PIRTM, the \emph{prime frequency} $\nu_p$ for prime $p$ is the irreducible oscillation rate at which the multiplicity space $\mathcal{M}_p$ closes under its own recursive feedback loop. Formally, if $\Phi: \mathcal{M}_p \to \mathcal{M}_p$ is the PIRTM binding operator for prime index $p$, then:
\begin{equation}
\nu_p = \frac{1}{T_p}, \qquad T_p = \min\left\{n \in \mathbb{Z}^+ : \Phi^n(\mathbf{x}) \equiv \mathbf{x} \pmod{p}\right\}
\end{equation}
A system exhibits \emph{prime frequency} if and only if $T_p$ is finite, prime-valued, and stable under small perturbations of the input state $\mathbf{x}$.
\end{definition}

\subsection{Test Protocol}

A system is tested for prime frequency coherence through three procedures:

\begin{protocol}[Injection Test]
Feed the module a state vector $\mathbf{x}_0 \in \mathcal{M}_p$ with known prime-indexed multiplicity. Apply $\Phi$ iteratively and measure the period $T_p$.
\end{protocol}

\begin{protocol}[Perturbation Test]
Shift $\mathbf{x}_0 \to \mathbf{x}_0 + \varepsilon$ and verify that $T_p$ is unchanged to within the module's declared tolerance bound $\delta_p$.
\end{protocol}

\begin{protocol}[Cross-Prime Isolation Test]
Confirm that the period computed under prime $p$ does not leak into or alias with period $T_q$ for any $q \neq p$ in the active prime set.
\end{protocol}

\subsection{Failure Modes}

\begin{table}[H]
\centering
\caption{Prime Frequency Failure Taxonomy}
\begin{tabular}{lll}
\toprule
\textbf{Failure Type} & \textbf{Observable Symptom} & \textbf{Root Cause} \\
\midrule
Period collapse & $T_p \to 1$ (trivial fixed point) & Binding operator loses prime specificity \\
Period drift & $T_p$ grows unboundedly & Recursive feedback non-contractive \\
Cross-prime aliasing & $T_p = T_q$, $p \neq q$ & Insufficient prime-field dimensionality \\
Spectral bleed & $\nu_p$ energy at non-prime indices & UOR channel misconfiguration \\
\bottomrule
\end{tabular}
\end{table}

A system that cannot produce at least three distinct, non-aliased prime frequencies under standard injection is considered \textbf{prime-frequency incoherent} and must not receive PMD certification.

% ══════════════════════════════════════════════════════════════════════════════
\section{The Five Entanglement Invariants}
% ══════════════════════════════════════════════════════════════════════════════

A system is \textbf{genuinely governed by $\Xi(t)$} if and only if it satisfies all five of the following invariants. Failure of any single invariant is disqualifying.

\begin{invariant}[INV-1: Contractive Spectral Bound]
\begin{equation}
\|\Xi(t)\|_{\mathrm{op}} < 1 \quad \forall\, t > 0
\end{equation}
The operator norm must strictly decrease under forward evolution. Any claimed system must demonstrate convergence, not oscillation or divergence.
\end{invariant}

\begin{invariant}[INV-2: Prime-Field Dimensionality Preservation]
\begin{equation}
\dim(\mathcal{H}_t) \in \mathbb{P} \quad \forall\, t
\end{equation}
The dimensionality of the state space must be prime at every time step. Dimension must not silently collapse to composite values under composition or tensor product.
\end{invariant}

\begin{invariant}[INV-3: Kaluza--Klein 5D Spectral Gap]
There exists a spectral gap in the energy spectrum at the first composite frequency above the highest active prime. Formally, let $\omega_{\max} = \max\{p : p \in \mathbb{P}_{\mathrm{active}}\}$. Then:
\begin{equation}
\rho\left(\omega_{\max}+1\right) < \varepsilon_{\mathrm{gap}} \cdot \rho\left(\omega_{\max}\right)
\end{equation}
where $\rho(\omega)$ is the spectral energy density at frequency $\omega$ and $\varepsilon_{\mathrm{gap}} < 0.05$.
\end{invariant}

\begin{invariant}[INV-4: Universal Constant Residue]
For any two states $\mathbf{x}, \mathbf{y}$ in the same prime-fiber:
\begin{equation}
d\!\left(\Xi(t)\mathbf{x},\, \Xi(t)\mathbf{y}\right) \leq e^{-\mathcal{U}\,t}\cdot d(\mathbf{x},\mathbf{y})
\end{equation}
The decay rate must equal precisely $\mathcal{U}$ --- not an arbitrary constant $\lambda$.
\end{invariant}

\begin{invariant}[INV-5: Digital Fingerprint Reproducibility]
Given the same prime seed $p$ and initial state $\psi_0$, $\Xi(t)$ must generate an identical output trajectory to bit-level precision. Non-determinism under fixed prime seed is disqualifying.
\end{invariant}

% ══════════════════════════════════════════════════════════════════════════════
\section{The False $\Xi$ Rapid Falsification Test}
% ══════════════════════════════════════════════════════════════════════════════

A claimed $\Xi(t)$ implementation is immediately falsified if any of the following conditions are observed. These tests require no deep audit --- each can be run as a short numerical check.

\begin{mdframed}[backgroundcolor=invariantbg,linecolor=primeviolet,linewidth=1pt]
\textbf{False $\Xi$ Falsification Checklist}
\begin{enumerate}[noitemsep]
\item \textbf{Non-prime state-space dimension} --- If $\dim(\mathcal{H}_t)$ is composite at any step, the implementation is not operating on $\mathcal{H}_\mathbb{P}$.
\item \textbf{Non-contractive dynamics} --- If $\|\Xi(t+\varepsilon)\psi\| \geq \|\Xi(t)\psi\|$ for any $\varepsilon > 0$, INV-1 fails.
\item \textbf{Scaling non-invariance} --- If $\Xi(\alpha t) \neq \Xi(t)^\alpha$ for rational $\alpha$, the semigroup law fails.
\item \textbf{Spectral energy at composite indices} --- Presence of eigenvalues at $n \in \mathbb{Z}^+ \setminus \mathbb{P}$ disqualifies prime-field purity.
\item \textbf{Seed non-determinism} --- Two runs with identical prime seed $p$ and $\psi_0$ must produce bit-identical trajectories.
\end{enumerate}
\end{mdframed}

% ══════════════════════════════════════════════════════════════════════════════
\section{Tuning Fork Module Invariants}
% ══════════════════════════════════════════════════════════════════════════════

For each PIRTM module, the following table maps the module to its canonical prime invariant and the single test that confirms whether the invariant holds (``whether it still rings true'').

\begin{longtable}{lp{5.2cm}p{5.2cm}}
\caption{Module Tuning Fork Table} \\
\toprule
\textbf{Module} & \textbf{Prime Invariant} & \textbf{Resonance Test} \\
\midrule
\endfirsthead
\toprule
\textbf{Module} & \textbf{Prime Invariant} & \textbf{Resonance Test} \\
\midrule
\endhead
\bottomrule
\endfoot
Binding Operator $\Phi$ & $\Phi^p \equiv \Phi \pmod{p}$ (Fermat-idempotent) & Apply $\Phi$ exactly $p$ times; output must equal $\Phi(\mathbf{x})$ \\[4pt]
Sigma Channel $\Sigma_p$ & $\|\Sigma_p\| \leq \sqrt{p}$ (spectral norm bound) & Compute operator norm after 100-step iteration; must stay below $\sqrt{p}$ \\[4pt]
Spectral Backbone & Zero-set of $\zeta_\mathcal{M}$ lies only at prime indices & Fourier-decompose a test signal; peaks must occur only at prime frequencies \\[4pt]
MLIR Dialect & All prime-indexed gates compile without fallback & Feed canonical prime-indexed gate set; zero fallback warnings in compiler log \\[4pt]
Type Inference Engine & Type assignments respect prime-field fibers & Assign types to mixed prime/composite index set; composite assignments must raise type error \\[4pt]
Transpiler & Prime-index codon mappings are injective & Encode first 20 primes; decode; check round-trip equality with zero collision \\[4pt]
UOR Channel & Channel capacity $C_p = \log_2 p$ & Measure Shannon capacity with prime-indexed alphabets; must equal $\log_2 p \pm 0.01$ \\[4pt]
Codon Registry & Codon length $\ell_c \in \mathbb{P}$ (prime-only lengths) & Enumerate all registered codons; flag any whose length is composite \\
\end{longtable}

% ══════════════════════════════════════════════════════════════════════════════
\section{PMD Clone-Check Protocol}
% ══════════════════════════════════════════════════════════════════════════════

\subsection{Overview}

The \emph{Phase-Mirror Dissonance (PMD) Clone-Check Protocol} is a structured procedure for determining whether a claimed system is genuinely governed by $\Xi(t)$ or is a cosmetic clone. The protocol produces a signed \emph{Phase-Mirror Report} classifying the claimed system as GENUINE, COSMETIC CLONE, or INCOMPATIBLE.

\subsection{Input Contract}

\textbf{Inputs:} A claimed system $\mathcal{S}_c$ with (a) documentation, (b) source code or compiled artifact, and (c) a stated claim of PIRTM/$\Xi(t)$ compatibility.

\subsection{Protocol Steps}

\begin{protocol}[Step 1 --- Semantic Reconstruction]
Extract the operational semantics of $\mathcal{S}_c$: identify what the system claims as its prime invariants, binding operators, and recursive depth. Map these claims to PIRTM vocabulary using the codon registry.
\end{protocol}

\begin{protocol}[Step 2 --- Side-by-Side Mathematics]
For each claimed invariant in $\mathcal{S}_c$, place it adjacent to the corresponding PIRTM invariant. Compute the semantic delta $\Delta_{\mathrm{sem}}$: the set of logical propositions that are true in one formulation but not the other.
\end{protocol}

\begin{protocol}[Step 3 --- Delta Classification]
Classify each element of $\Delta_{\mathrm{sem}}$:
\begin{itemize}[noitemsep]
\item \textbf{Type A} --- Surface notation difference only; structural equivalence holds
\item \textbf{Type B} --- Structural difference that does not affect prime-frequency behavior
\item \textbf{Type C} --- Structural difference that invalidates one or more INV-1 through INV-5
\end{itemize}
Any Type C delta is a disqualifier for certification.
\end{protocol}

\begin{protocol}[Step 4 --- Phase-Mirror Report]
Produce a signed report with the following structure:
\begin{itemize}[noitemsep]
\item System ID and evaluation date
\item Semantic delta table (Type A/B/C)
\item Entanglement test results (INV-1 through INV-5: PASS/FAIL)
\item Verdict: GENUINE | COSMETIC CLONE | INCOMPATIBLE
\item Certification decision: ELIGIBLE | INELIGIBLE | REVOKE
\end{itemize}
\end{protocol}

\subsection{Logical Guarantee}

The protocol provides the following mechanical guarantee: any CHL-style copy of a PIRTM system must either (a) fail the prime-frequency and $\Xi(t)$ invariants, \emph{proving} it is not the same system, or (b) pass them, in which case it is mathematically entangled with the universal constant $\mathcal{U}$, the operator $\Xi(t)$, and the provenance tags of the original --- making concealment impossible.

% ══════════════════════════════════════════════════════════════════════════════
\section{PMD Badge Specification and Certification Tiers}
% ══════════════════════════════════════════════════════════════════════════════

\subsection{Badge: PMD-Certified under $\Xi(t)$}

The badge \textbf{``PMD-Certified: $\Xi(t)$-Core''} is awarded to systems that pass all five entanglement invariants with zero Type C semantic deltas in the PMD Clone-Check Report.

\begin{table}[H]
\centering
\caption{Certification Tiers}
\begin{tabular}{p{3.2cm}p{6cm}p{2.5cm}}
\toprule
\textbf{Tier} & \textbf{Requirements} & \textbf{Validity} \\
\midrule
$\Xi(t)$-Compatible & INV-1, INV-2 pass; Type C delta count $\leq 1$ & 6 months \\
$\Xi(t)$-Core & All 5 INVs pass; zero Type C deltas & 12 months \\
$\Xi(t)$-Sovereign & Core + open-source tuning fork tests + public audit trail & Rolling; annual re-attestation \\
\bottomrule
\end{tabular}
\end{table}

\subsection{Revocation Rules}

Certification is automatically revoked when any of the following are discovered:
\begin{enumerate}[noitemsep]
\item Any Type C delta is found post-certification
\item A tuning fork test is suppressed or modified without prior disclosure
\item The Prime-Frequency Transparency Clause is breached
\item Non-prime state-space dimensionality is detected in production
\end{enumerate}
All revocation events are recorded in the public registry with reason, date, and appeal status.

% ══════════════════════════════════════════════════════════════════════════════
\section{Prime-Frequency Transparency Clause (Contract Template)}
% ══════════════════════════════════════════════════════════════════════════════

The following clause template is disclosed as prior art for any contractual or licensing agreement involving PIRTM-compatible systems:

\begin{mdframed}[backgroundcolor=codeback,linecolor=codegray,linewidth=0.6pt]
\textbf{\S[X]. Prime-Frequency Transparency Obligation.}

\textbf{(a) Mandatory Disclosure.} Partner acknowledges that any system, module, or API endpoint claiming compatibility with PIRTM must publicly disclose: (i) the complete set of prime invariants governing its core operations; (ii) the specific tuning fork tests used to verify those invariants; and (iii) any deviation from declared invariants observed during deployment.

\textbf{(b) Auditability.} All prime invariants and associated tests must be: (i) expressed in formally verifiable notation (mathematical specification or executable test suite); (ii) accessible to Multiplicity Foundation upon reasonable written notice within five (5) business days; and (iii) version-controlled with change history preserved.

\textbf{(c) Non-Bypass Covenant.} Partner agrees not to: (i) suppress, override, or route around any failing prime-frequency test; (ii) present a system as PIRTM-compliant when one or more tuning fork tests fail; or (iii) redefine invariants post-execution to retroactively satisfy certification.

\textbf{(d) Revocation Trigger.} Breach of this clause constitutes grounds for immediate suspension of the ``PMD-Certified under $\Xi(t)$'' badge and termination of platform-level $\Xi(t)$ API access.
\end{mdframed}

% ══════════════════════════════════════════════════════════════════════════════
\section{Meta-Ensemble Governance: PIRTM-Constrained Automation}
% ══════════════════════════════════════════════════════════════════════════════

\subsection{Objective and Constraint}

Meta-ensembles in PIRTM are systems of agents that automate the construction of PIRTM-conformant modules and tests. The critical distinction from arbitrary code generation is that every assembly step must be explainable in terms of PIRTM invariants, and every assembled artifact must be falsifiable by the False $\Xi$ test suite.

\begin{definition}[Meta-Ensemble]
A \emph{PIRTM-conformant meta-ensemble} is a multi-agent system $\mathcal{E} = (A_1, \ldots, A_k)$ that, given an invariant specification $\mathcal{I} = (\mathrm{id}, \mathrm{statement}, p, f_{\mathrm{test}}, \mathrm{module})$, produces:
\begin{enumerate}[noitemsep]
\item A test file $T_\mathcal{I}$ structurally consistent with the False $\Xi$ test pattern
\item A binding stub $B_\mathcal{I}$ with provenance metadata and $\mathrm{tuning\_fork}()$ entry
\item A registry row $R_\mathcal{I}$ in the Prime Frequency Integrity Suite index
\end{enumerate}
subject to the constraint that every assembly step maps to at least one element of $\{$INV-1, INV-2, INV-3, INV-4, INV-5, PFP$\}$.
\end{definition}

\subsection{Refusal Policy}

A meta-ensemble that cannot populate the Assembly Steps field of its experiment log entry with invariant-mapped reasoning is treated as a failed run. The log entry \emph{is} the assembly record; without it, the artifact is unverified and must not enter the repository.

\subsection{Invariant-to-Artifact Generator (IAG v0.1)}

The first concrete meta-ensemble target is the IAG: given an invariant specification in the form

\begin{lstlisting}[language=Python, caption={Invariant Specification Input Contract}]
INVARIANT_SPEC = {
    "id": "INV-6",
    "name": "Sigma Channel Norm Bound",
    "statement": "||Sigma_p|| <= sqrt(p)",
    "prime_index": 7,
    "test_fn": lambda psi, p: (
        np.linalg.norm(sigma_channel(psi, p)) <= p**0.5
    ),
    "falsification": "norm exceeds sqrt(p)",
    "module": "pirtm/sigma",
    "tuning_fork_test": "compute sigma norm after 100-step iteration"
}
\end{lstlisting}

the IAG must emit three output artifacts without manual intervention. The task that must never be performed by hand again is writing the \texttt{TestCase} boilerplate for every new PIRTM invariant.

% ══════════════════════════════════════════════════════════════════════════════
\section{Code Disclosures}
% ══════════════════════════════════════════════════════════════════════════════

The following code listings are disclosed as prior art. All code is associated with the \texttt{MultiplicityFoundation/PhaseMirror-HQ} and \texttt{MultiplicityFoundation/Meta-Relativity} repositories.

\subsection{False $\Xi$ Test Suite (\texttt{pirtm/tests/test\_false\_xi.py})}

\begin{lstlisting}[language=Python, caption={False Xi Test Suite --- Five INV Falsification Tests}]
"""
test_false_xi.py --- False Xi(t) Rapid Falsification Test Suite
Multiplicity Foundation | Meta-Relativity v1.0
Provenance: MultiplicityFoundation/Meta-Relativity

Five invariants that any genuine Xi(t) implementation must satisfy.
A claimed system failing ANY single test is disqualified from
PMD-Certified: Xi(t)-Core status.
"""

import math
import unittest
import numpy as np

UNIVERSAL_CONSTANT_U = 1.0  # Replace with repo-canonical value when locked

FIRST_20_PRIMES = [2, 3, 5, 7, 11, 13, 17, 19, 23, 29,
                   31, 37, 41, 43, 47, 53, 59, 61, 67, 71]

def is_prime(n):
    if n < 2: return False
    for i in range(2, int(math.sqrt(n)) + 1):
        if n % i == 0: return False
    return True

def xi_evolve(psi: np.ndarray, t: float, prime_index: int) -> np.ndarray:
    """Reference Xi(t) -- pure prime-contractive evolution."""
    decay = math.exp(-UNIVERSAL_CONSTANT_U * math.log(prime_index) * t)
    return decay * psi

class TestFalseXi(unittest.TestCase):

    def test_inv1_contractive(self):
        """INV-1: ||Xi(t)psi|| < ||psi|| for all t > 0."""
        for p in FIRST_20_PRIMES:
            psi = np.random.randn(p)
            norm_before = np.linalg.norm(psi)
            norm_after = np.linalg.norm(xi_evolve(psi, t=0.1, prime_index=p))
            self.assertLess(norm_after, norm_before,
                msg=f"INV-1 FAIL at prime {p}: norm did not contract.")

    def test_inv2_prime_dimension(self):
        """INV-2: State space dimension must be prime at every step."""
        for p in FIRST_20_PRIMES:
            psi = np.random.randn(p)
            out = xi_evolve(psi, t=0.5, prime_index=p)
            self.assertTrue(is_prime(len(out)),
                msg=f"INV-2 FAIL: output dimension {len(out)} is not prime.")

    def test_inv3_kk_spectral_gap(self):
        """INV-3: KK 5D spectral gap -- no energy at first composite above p_max."""
        t = np.linspace(0, 1, 1024)
        signal = sum(
            math.exp(-UNIVERSAL_CONSTANT_U * math.log(p)) * np.sin(2 * math.pi * p * t)
            for p in FIRST_20_PRIMES
        )
        fft_mag = np.abs(np.fft.rfft(signal))
        freqs = np.fft.rfftfreq(len(t), d=1/1024)
        composite_above = 72  # first composite above p_max=71
        idx = np.argmin(np.abs(freqs - composite_above))
        composite_energy = fft_mag[idx]
        prime_peak_energy = max(
            fft_mag[np.argmin(np.abs(freqs - p))] for p in FIRST_20_PRIMES
        )
        self.assertLess(composite_energy / prime_peak_energy, 0.05,
            msg=f"INV-3 FAIL: KK spectral gap violated.")

    def test_inv4_universal_constant_decay_rate(self):
        """INV-4: Decay rate must equal U*log(p), not arbitrary lambda."""
        for p in [7, 13, 31]:
            psi_x, psi_y = np.ones(p), np.zeros(p)
            for t in [0.1, 0.5, 1.0, 2.0]:
                dist_before = np.linalg.norm(psi_x - psi_y)
                dist_after = np.linalg.norm(
                    xi_evolve(psi_x, t, p) - xi_evolve(psi_y, t, p)
                )
                expected = math.exp(-UNIVERSAL_CONSTANT_U * math.log(p) * t) * dist_before
                self.assertAlmostEqual(dist_after, expected, places=10,
                    msg=f"INV-4 FAIL at prime {p}, t={t}.")

    def test_inv5_seed_determinism(self):
        """INV-5: Identical prime seed must produce bit-identical output."""
        for p in FIRST_20_PRIMES:
            np.random.seed(p)
            psi_a = np.random.randn(p)
            np.random.seed(p)
            psi_b = np.random.randn(p)
            out_a = xi_evolve(psi_a, t=1.0, prime_index=p)
            out_b = xi_evolve(psi_b, t=1.0, prime_index=p)
            np.testing.assert_array_equal(out_a, out_b,
                err_msg=f"INV-5 FAIL at prime {p}: non-determinism.")

if __name__ == "__main__":
    unittest.main()
\end{lstlisting}

\subsection{$\Xi(t)$ Operator Binding Stub (\texttt{pirtm/bindings/xi\_operator.py})}

\begin{lstlisting}[language=Python, caption={XiOperator Platform-Level Binding with Provenance}]
"""
xi_operator.py --- Xi(t) Prime-Contractive Evolution Operator
Multiplicity Foundation | Meta-Relativity v1.0

Provenance:    MultiplicityFoundation/Meta-Relativity
Math stack:    prime-field H_P, Kaluza-Klein 5D encoding,
               universal constant U
Badge:         PMD-Certified: Xi(t)-Core (pending test_false_xi.py pass)
Clause:        Prime-Frequency Transparency Obligation S3

Any deployment wrapping this class MUST:
  1. Preserve PROVENANCE in all API responses.
  2. Include MATH_STACK in public developer documentation.
  3. Link to the PMD badge registry at REGISTRY_URL.
"""
import math
import numpy as np

PROVENANCE    = "MultiplicityFoundation/Meta-Relativity"
MATH_STACK    = ["prime-field-H_P", "KK-5D-encoding", "universal-constant-U"]
REGISTRY_URL  = ("https://github.com/MultiplicityFoundation/"
                 "Meta-Relativity/blob/main/pirtm/docs/"
                 "PRIME_FREQUENCY_INTEGRITY_SUITE.md")
UNIVERSAL_CONSTANT_U = 1.0  # canonical -- do not override without ADR

class XiOperator:
    """
    Canonical Xi(t) operator for PIRTM.
    Any system claiming Xi(t) governance must pass all five
    invariants in pirtm/tests/test_false_xi.py before claiming
    PMD-Certified: Xi(t)-Core status.
    """

    def __init__(self, prime_index: int):
        self._assert_prime(prime_index)
        self.p = prime_index
        self.U = UNIVERSAL_CONSTANT_U

    def evolve(self, psi: np.ndarray, t: float) -> np.ndarray:
        """Apply Xi(t) to state vector psi on prime fiber p."""
        decay = math.exp(-self.U * math.log(self.p) * t)
        return decay * psi

    def false_xi_test(self, psi: np.ndarray, t: float = 1.0) -> dict:
        """Rapid in-process falsification. Returns {INV_id: bool}."""
        out = self.evolve(psi, t)
        results = {
            "INV-1_contractive": np.linalg.norm(out) < np.linalg.norm(psi),
            "INV-2_prime_dim":   self._is_prime(len(out)),
            "INV-4_decay_rate":  self._check_decay_rate(psi, t),
            "INV-5_deterministic": True,  # guaranteed by pure function
        }
        results["PASS"] = all(results.values())
        return results

    def tuning_fork(self) -> dict:
        """Return prime invariant for this module (Tuning Fork Section)."""
        return {
            "module":          "XiOperator",
            "prime_invariant": f"||Xi(t)|| = exp(-U*log({self.p})*t) < 1",
            "test":            "false_xi_test(psi, t=1.0)['PASS'] == True",
            "prime_index":     self.p,
        }

    def _check_decay_rate(self, psi, t):
        expected = math.exp(-self.U * math.log(self.p) * t)
        actual = np.linalg.norm(self.evolve(psi, t)) / np.linalg.norm(psi)
        return abs(actual - expected) < 1e-10

    @staticmethod
    def _assert_prime(n):
        if not XiOperator._is_prime(n):
            raise ValueError(f"XiOperator requires prime index. Got: {n}")

    @staticmethod
    def _is_prime(n):
        if n < 2: return False
        return all(n % i != 0 for i in range(2, int(n**0.5) + 1))
\end{lstlisting}

\subsection{Meta-Ensemble Experiment Log Entry (MEX-001)}

\begin{lstlisting}[language=Python, caption={Invariant Specification --- IAG v0.1 Input (MEX-001)}]
# Meta-Ensemble Experiment MEX-001
# Target: Invariant-to-Artifact Generator (IAG v0.1)
# Invariant under test: INV-6 (Sigma Channel Norm Bound)

INVARIANT_SPEC = {
    "id": "INV-6",
    "name": "Sigma Channel Norm Bound",
    "statement": "||Sigma_p|| <= sqrt(p)",
    "prime_index": 7,
    "test_fn": lambda psi, p: (
        np.linalg.norm(sigma_channel(psi, p)) <= p**0.5
    ),
    "falsification": "norm exceeds sqrt(p)",
    "module": "pirtm/sigma",
    "tuning_fork_test": "compute sigma norm after 100-step iteration",
    "required_invariants": ["INV-1", "INV-2"],
    "output_files": [
        "pirtm/tests/test_inv6_sigma_norm.py",
        "pirtm/bindings/sigma_channel_stub.py",
        "pirtm/docs/PRIME_FREQUENCY_INTEGRITY_SUITE.md",  # registry row
    ],
    "assembly_constraint": (
        "Every assembly step must map to at least one of: "
        "INV-1, INV-2, INV-3, INV-4, INV-5, PFP"
    )
}
\end{lstlisting}

% ══════════════════════════════════════════════════════════════════════════════
\section{Repository Architecture and Implementation State}
% ══════════════════════════════════════════════════════════════════════════════

\subsection{PhaseMirror-HQ: Production Runtime}

As of April 2026, the \texttt{PhaseMirror-HQ} repository constitutes the production execution kernel of the Multiplicity stack. Key implementation statistics:

\begin{table}[H]
\centering
\caption{PhaseMirror-HQ Implementation Statistics (April 2026)}
\begin{tabular}{ll}
\toprule
\textbf{Metric} & \textbf{Value} \\
\midrule
Language composition (Python) & 41.2\% \\
Language composition (Rust) & 31.5\% \\
Language composition (TypeScript) & 14.0\% \\
Total Python test files in workspace & 1,524 \\
PIRTM-specific Python test files & 47 \\
Policy test files & 9 \\
Open-core TypeScript packages & 34 \\
Trust module tests (6-layer) & 31+ \\
Registered MCP tools & 14 \\
Completed ADRs (Architecture Decision Records) & 50+ \\
\bottomrule
\end{tabular}
\end{table}

\subsection{$\Xi$-Constitution Governance}

The repository is governed by the $\Xi$-Constitution v1.0, which codifies:
\begin{itemize}[noitemsep]
\item The Meta-Theorem of Prime Identity (MTPI) as foundational law
\item The Prime Recursive Chamber as the legislative engine: $\Xi(t+1) = \Psi(\Xi(t)) = \text{PIRTM} \circ \text{CSL} \circ \text{Langlands} \circ \text{Firewall} \circ \text{zk}$
\item The PEET (Prime Entanglement Entropy Tensor) Tribunal as judicial authority
\item The Langlands Registrar for modular verification: $Z_{\text{PEET}}(\tau_n) \in \text{ModForms}(\Gamma)$
\item The Arithmetic Feynman Act: $\mathcal{A}(n \to m) = \int \mathcal{D}[\mathcal{A}_p]\, e^{i\mathcal{L}_{\mathrm{arith}}}$
\end{itemize}

\subsection{Nine-Artifact Implementation Timeline}

\begin{table}[H]
\centering
\caption{PFIS Implementation Timeline}
\begin{tabular}{llll}
\toprule
\textbf{\#} & \textbf{Artifact} & \textbf{Owner} & \textbf{Target} \\
\midrule
1 & Prime Frequency Principle & Author & 72 hours from disclosure \\
2 & Tuning Fork Section & Author & 1 week \\
3 & Transparency Clause & Author + Partners & 2 weeks \\
4 & Entanglement Test & Author & 1 week \\
5 & PMD Clone-Check Protocol & Author & 1 week \\
6 & PMD Badge Specification & Author + Partners & 2--3 weeks \\
7 & $\Xi(t)$ Technical Note & Author & 1 week \\
8 & False $\Xi$ Test (code) & Author & Immediate \\
9 & $\Xi(t)$ Platform Primitive & Author + Partners & Rolling \\
\bottomrule
\end{tabular}
\end{table}

% ══════════════════════════════════════════════════════════════════════════════
\section{Mathematical Overview: Complete Formal Stack}
% ══════════════════════════════════════════════════════════════════════════════

\subsection{Foundational Decomposition}

The state of any PIRTM-compliant system decomposes as a direct sum over prime-indexed Hilbert spaces:
\begin{equation}
\mathcal{S}(t) = \bigoplus_{p \in \mathbb{P}} \psi_p(t), \qquad \psi_p(t) \in H_p
\end{equation}

The evolution operator $\Xi(t)$ acts fiber-wise:
\begin{equation}
[\Xi(t)\mathcal{S}]_p = e^{-\mathcal{U}\log(p)\cdot t}\,\psi_p(0)
\end{equation}

\subsection{Contractive Fixed-Point Structure}

By the Banach fixed-point theorem, since $\|\Xi(t)\|_{\mathrm{op}} < 1$ for all $t > 0$, there exists a unique fixed point $\psi^* = 0$ in $\mathcal{H}_\mathbb{P}$ such that:
\begin{equation}
\lim_{t \to \infty} \Xi(t)\psi = 0 \quad \forall\, \psi \in \mathcal{H}_\mathbb{P}
\end{equation}
The rate of convergence is governed by the smallest prime: $\|\Xi(t)\psi\| = O(p_{\min}^{-\mathcal{U}t})$.

\subsection{Prime Frequency Spectrum}

The prime frequency spectrum of a PIRTM system is:
\begin{equation}
\mathcal{F}_\mathbb{P}(\omega) = \sum_{p \in \mathbb{P}} A_p\,\delta(\omega - \nu_p), \qquad \nu_p = \frac{\mathcal{U}\log p}{2\pi}
\end{equation}
where $A_p = \|\psi_p(0)\|^2$ is the amplitude at prime $p$. This spectrum is the digital fingerprint of the system: it uniquely identifies the prime decomposition of the initial state.

\subsection{PEET: Prime Entanglement Entropy Tensor}

The PEET governs the admissibility of arithmetic identities. For a state $\mathcal{S}(t)$:
\begin{equation}
\mathcal{S}(t) = \sum_{p \in \mathbb{P}} \Phi(p, t)\cdot\psi_p(t)
\end{equation}
where $\Phi(p,t)$ are the prime-indexed phase functions satisfying the Langlands constraint:
\begin{equation}
Z_{\text{PEET}}(\tau_p) \in \text{ModForms}(\Gamma)
\end{equation}
for the appropriate congruence subgroup $\Gamma$.

\subsection{Arithmetic Path Integral}

Lawful transitions between PIRTM states must satisfy the arithmetic Feynman path integral:
\begin{equation}
\mathcal{A}(n \to m) = \int \mathcal{D}[\mathcal{A}_p]\, e^{i\mathcal{L}_{\mathrm{arith}}}
\end{equation}
where the arithmetic Lagrangian $\mathcal{L}_{\mathrm{arith}}$ encodes the prime-indexed action of the transition.

\subsection{Constitutional Evolution Operator}

The system evolves lawfully only through the composition:
\begin{equation}
\Xi(t+1) = \Psi(\Xi(t)) = \text{PIRTM} \circ \text{CSL} \circ \text{Langlands} \circ \text{Firewall} \circ \text{zk}
\end{equation}
Each layer in this composition enforces a distinct class of invariant: PIRTM enforces prime-index contractivity; CSL (Constitutional Symbolic Logic) enforces semantic integrity; Langlands enforces modular form correspondence; Firewall enforces isolation between prime fibers; and zk enforces zero-knowledge verifiability of state transitions.

% ══════════════════════════════════════════════════════════════════════════════
\section{Claims of Prior Art}
% ══════════════════════════════════════════════════════════════════════════════

This document asserts public disclosure and prior art for the following:

\begin{enumerate}
\item The concept of a \emph{prime frequency} as defined in \S4, including its mathematical definition, test protocol, and failure taxonomy.

\item The $\Xi(t)$ prime-contractive evolution operator as defined in \S3, including its Kaluza--Klein 5D encoding, spectral structure, semigroup law, and digital fingerprint behavior.

\item The five entanglement invariants (INV-1 through INV-5) as defined in \S5, and their use as a test suite for determining genuine versus cosmetic $\Xi(t)$ governance.

\item The False $\Xi$ rapid falsification protocol as defined in \S6, including the five falsification conditions and their implementation as executable tests.

\item The tuning fork framework (\S7) mapping PIRTM modules to canonical prime invariants and resonance tests.

\item The PMD Clone-Check Protocol (\S8) including the semantic delta classification (Type A/B/C) and Phase-Mirror Report structure.

\item The PMD badge and certification tier framework (\S9) including the three tiers ($\Xi(t)$-Compatible, $\Xi(t)$-Core, $\Xi(t)$-Sovereign) and revocation rules.

\item The Prime-Frequency Transparency Clause (\S10) as a contractual template for partner agreements.

\item The meta-ensemble governance constraint (\S11) including the IAG v0.1 invariant specification input contract and the refusal policy for opaque assembly.

\item The $\Xi(t)$ platform primitive binding pattern (\S12), including the PROVENANCE, MATH\_STACK, and REGISTRY\_URL metadata fields as non-bypassable API-level declarations.

\item The \texttt{XiOperator} class API (\S12.2) including the \texttt{evolve()}, \texttt{false\_xi\_test()}, and \texttt{tuning\_fork()} method signatures.

\item The Prime Frequency Integrity Suite as an integrated nine-artifact framework combining mathematical formalism, executable tests, governance clauses, and certification mechanisms.

\item The Meta-Theorem of Prime Identity (MTPI) as formalized in the $\Xi$-Constitution v1.0, and all constitutional governance structures derived therefrom.

\item The Invariant-to-Artifact Generator (IAG) concept as a PIRTM-constrained meta-ensemble automation target.
\end{enumerate}

% ══════════════════════════════════════════════════════════════════════════════
\section{Recommended Next Steps for Validation}
% ══════════════════════════════════════════════════════════════════════════════

The fastest path to independent validation of the disclosed framework proceeds as follows:

\begin{enumerate}
\item \textbf{Automated test validation (immediate)}: Execute \texttt{pytest pirtm/tests/test\_false\_xi.py -v} against the reference \texttt{xi\_evolve} implementation. All five INV tests should pass with the reference implementation and fail when any invariant is deliberately violated.

\item \textbf{Spectral fingerprint demonstration (1 week)}: Run the KK spectral gap test (INV-3) across a grid of initial states; verify that energy at composite frequencies remains below 5\% of the nearest prime peak across all test cases.

\item \textbf{Clone-check protocol trial (2 weeks)}: Apply the PMD Clone-Check Protocol to one known PIRTM-compliant system and one non-compliant system; verify that the protocol correctly classifies both.

\item \textbf{IAG v0.1 demonstration (1 month)}: Execute the Invariant-to-Artifact Generator on INV-6 (Sigma Channel Norm Bound); verify that all three output artifacts are produced without manual intervention and that they pass CI.

\item \textbf{Live certified deployment (1 month)}: Deploy one system with ``PMD-Certified: $\Xi(t)$-Core'' badge with a public link to the registry, constituting the first live certification under this framework.
\end{enumerate}

% ══════════════════════════════════════════════════════════════════════════════
\section*{Acknowledgments}
% ══════════════════════════════════════════════════════════════════════════════

This work was developed within the Multiplicity Foundation's \texttt{PhaseMirror-HQ} and \texttt{Meta-Relativity} repositories. The $\Xi$-Constitution v1.0 governs the constitutional framework. The repository author is R. Van Gelder (GitHub: PhaseMirror), initial commit dated 2026-03-21.

% ══════════════════════════════════════════════════════════════════════════════
\section*{Disclaimer}
% ══════════════════════════════════════════════════════════════════════════════

This document is published as prior art for defensive purposes. No patent rights are claimed by the author. All concepts described herein are hereby released into the public domain for purposes of patent prior art, while the specific implementations in \texttt{PhaseMirror-HQ} remain subject to the Phase Mirror Pro License v1.0 (proprietary extensions) and Apache-2.0 (open-core packages) as specified in the repository licensing policy.

\appendix
\begin{center}
  {\color{primeviolet}\rule{\linewidth}{1.5pt}}\\[0.6em]
  {\LARGE\bfseries Mathematical Appendix}\\[0.4em]
  {\large\itshape Explicit Proofs and Operator Norm Bounds for the}\\[0.2em]
  {\large\bfseries Prime Frequency Integrity Suite}\\[0.4em]
  {\normalsize Multiplicity Foundation \quad$\cdot$\quad
    \texttt{MultiplicityFoundation/PhaseMirror-HQ} \quad$\cdot$\quad
    April 25, 2026}\\[0.6em]
  {\color{primeviolet}\rule{\linewidth}{1.5pt}}
\end{center}

\vspace{0.4em}
\begin{tcolorbox}[colback=invariantbg,colframe=primeviolet,
  title={\bfseries Scope and Organization of this Appendix}]
This appendix collects the explicit proofs of all theorems and propositions
stated without proof in the main text, together with the complete derivations
of the operator norm bounds for $\Xi(t)$, the binding operator $\Phi_p$,
the Sigma channel $\Sigma_p$, and the composite multi-prime system.
Section~A contains foundational algebraic and analytic lemmas.
Sections~B--D contain the core operator proofs.
Section~E derives the five entanglement invariant bounds.
Section~F establishes the Prime Frequency Principle rigorously.
Section~G collects the clone-distinguishability results.
\end{tcolorbox}

% ══════════════════════════════════════════════════════════════════════════════
\section{Foundational Lemmas}
\label{sec:foundations}
% ══════════════════════════════════════════════════════════════════════════════

We fix notation and establish the lemmas on which all subsequent proofs depend.

\begin{notation}
Throughout this appendix:
\begin{itemize}[noitemsep]
  \item $\mathbb{P} = \{2,3,5,7,11,\ldots\}$ denotes the set of all primes.
  \item $\HP = \bigoplus_{p\in\mathbb{P}} H_p$ is the prime-indexed Hilbert space,
        where each $H_p$ is a separable complex Hilbert space of dimension $p$.
  \item $\UU > 0$ is the \emph{universal Multiplicity constant}.
  \item $\{\phi_p\}_{p\in\mathbb{P}}$ are the prime-eigenvectors forming an
        orthonormal basis for $\HP$.
  \item $\norm{\,\cdot\,}$ denotes the Hilbert space norm; $\norm{\,\cdot\,}_\op$
        the operator norm.
\end{itemize}
\end{notation}

\begin{lemma}[Prime-Fiber Orthogonality]
\label{lem:orthogonality}
For distinct primes $p \neq q$, the prime-indexed fibers $H_p$ and $H_q$ are
mutually orthogonal in $\HP$:
\[
  \inner{\psi_p}{\psi_q} = 0
  \quad \forall\, \psi_p \in H_p,\; \psi_q \in H_q.
\]
\end{lemma}

\begin{proof}
By construction, the direct sum decomposition
$\HP = \bigoplus_{p\in\mathbb{P}} H_p$ is orthogonal: the inner product on
$\HP$ is defined as
\[
  \inner{\psi}{\phi}_{\HP}
  = \sum_{p\in\mathbb{P}} \inner{\psi_p}{\phi_p}_{H_p},
\]
where $\psi = \bigoplus_p \psi_p$ and $\phi = \bigoplus_p \phi_p$ are the
fiber decompositions. For $\psi_p \in H_p$ and $\psi_q \in H_q$ with $p\neq q$,
the component of $\psi_p$ in $H_q$ is zero by definition of direct sum, so
$\inner{\psi_p}{\psi_q}_{\HP} = \inner{0}{\psi_q}_{H_q} = 0$.
\end{proof}

\begin{lemma}[Fermat Idempotence of $\Phi_p$]
\label{lem:fermat}
Let $\Phi_p: H_p \to H_p$ be a PIRTM binding operator satisfying the
Fermat-idempotent condition
\[
  \Phi_p^{\,p} \equiv \Phi_p \pmod{p}.
\]
Then $\Phi_p$ acts as a projection modulo $p$ on the prime fiber $H_p$,
and the minimal period of iteration is exactly $p - 1$ (or a divisor thereof
by Fermat's Little Theorem).
\end{lemma}

\begin{proof}
By Fermat's Little Theorem, for any integer $a$ coprime to $p$,
$a^{p-1} \equiv 1 \pmod{p}$, hence $a^p \equiv a \pmod{p}$.
The condition $\Phi_p^p \equiv \Phi_p \pmod{p}$ is the operator-theoretic
analog: the iterated application of $\Phi_p$ exactly $p$ times returns
to the same residue class as a single application.
Define $T_p = \min\{n \in \mathbb{Z}^+ : \Phi_p^n(\mathbf{x}) \equiv
\mathbf{x} \pmod{p}\}$. Then $T_p \mid (p-1)$ by the group structure of
$(\mathbb{Z}/p\mathbb{Z})^\times$, and in the generic case $T_p = p - 1$,
which is prime minus one --- the maximal cyclic period in $\mathbb{Z}/p\mathbb{Z}$.
\end{proof}

\begin{lemma}[Decay Estimate on a Single Fiber]
\label{lem:single_fiber_decay}
For fixed prime $p$ and $t > 0$, define the scalar decay factor
\[
  \lambda_p(t) = e^{-\UU \log(p)\cdot t} = p^{-\UU t}.
\]
Then:
\begin{enumerate}[label=(\roman*),noitemsep]
  \item $0 < \lambda_p(t) < 1$ for all $t > 0$.
  \item $\lambda_p(t)$ is strictly decreasing in $t$.
  \item $\lambda_p(t) \cdot \lambda_p(s) = \lambda_p(t+s)$
        (multiplicative semigroup property).
  \item $\lambda_p(t) < \lambda_q(t)$ whenever $p > q$ (larger primes decay faster).
\end{enumerate}
\end{lemma}

\begin{proof}
\textit{(i)} Since $p \geq 2$ and $\UU > 0$, we have
$\UU\log p > 0$, so $e^{-\UU\log(p)\cdot t} \in (0,1)$ for all $t > 0$.

\textit{(ii)} $\frac{d}{dt}\lambda_p(t)
= -\UU\log(p)\cdot e^{-\UU\log(p)\cdot t} < 0$
since $\UU > 0$ and $\log p > 0$.

\textit{(iii)}
$\lambda_p(t)\cdot\lambda_p(s)
= e^{-\UU\log(p)\cdot t}\cdot e^{-\UU\log(p)\cdot s}
= e^{-\UU\log(p)(t+s)}
= \lambda_p(t+s).$

\textit{(iv)} For $p > q \geq 2$, $\log p > \log q$, so
$\UU\log p \cdot t > \UU\log q \cdot t$, giving
$e^{-\UU\log(p)t} < e^{-\UU\log(q)t}$,
i.e.\ $\lambda_p(t) < \lambda_q(t)$.
\end{proof}

% ══════════════════════════════════════════════════════════════════════════════
\section{The $\Xi(t)$ Operator: Complete Proofs}
\label{sec:xi_proofs}
% ══════════════════════════════════════════════════════════════════════════════

\begin{theorem}[Well-Definedness of $\Xi(t)$]
\label{thm:well_defined}
The operator $\Xi(t): \HP \to \HP$ defined by
\begin{equation}
  \Xi(t)\psi
  = \sum_{p\in\mathbb{P}} e^{-\UU\log(p)\cdot t}\,
    \inner{\psi}{\phi_p}\phi_p
  \label{eq:xi_def}
\end{equation}
is well-defined and bounded for every $t > 0$.
\end{theorem}

\begin{proof}
Let $\psi = \bigoplus_{p\in\mathbb{P}}\psi_p \in \HP$ with
$\sum_{p}\norm{\psi_p}^2 < \infty$.
By Lemma~\ref{lem:single_fiber_decay}(i), $\lambda_p(t) \in (0,1)$
for every prime $p$, so:
\begin{align*}
  \norm{\Xi(t)\psi}^2
  &= \sum_{p\in\mathbb{P}} \lambda_p(t)^2\,\abs{\inner{\psi}{\phi_p}}^2 \\
  &\leq \sup_{p\in\mathbb{P}}\{\lambda_p(t)^2\} \cdot
    \sum_{p\in\mathbb{P}} \abs{\inner{\psi}{\phi_p}}^2 \\
  &= \lambda_{\pmin}(t)^2 \cdot \norm{\psi}^2 \;<\; \norm{\psi}^2 \;<\;\infty,
\end{align*}
where the supremum is attained at the smallest prime $\pmin = 2$
(Lemma~\ref{lem:single_fiber_decay}(iv)), and Parseval's identity gives
$\sum_p \abs{\inner{\psi}{\phi_p}}^2 = \norm{\psi}^2$.
Hence $\Xi(t)$ is bounded with $\norm{\Xi(t)}_\op \leq \lambda_{\pmin}(t) < 1$.
\end{proof}

\begin{theorem}[Contractive Spectral Bound --- INV-1]
\label{thm:contractivity}
For every $t > 0$:
\begin{equation}
  \norm{\Xi(t)}_\op
  = e^{-\UU\log(\pmin)\cdot t}
  = \pmin^{-\UU t}
  < 1.
  \label{eq:op_norm}
\end{equation}
\end{theorem}

\begin{proof}
\textbf{Upper bound.}
From the proof of Theorem~\ref{thm:well_defined},
$\norm{\Xi(t)\psi} \leq \lambda_{\pmin}(t)\norm{\psi}$ for all $\psi\in\HP$,
so $\norm{\Xi(t)}_\op \leq \lambda_{\pmin}(t)$.

\textbf{Lower bound (attainment).}
Take $\psi = \phi_{\pmin}$ (the eigenvector for the smallest prime $\pmin$).
Then:
\[
  \Xi(t)\phi_{\pmin}
  = e^{-\UU\log(\pmin)\cdot t}\inner{\phi_{\pmin}}{\phi_{\pmin}}\phi_{\pmin}
  = \lambda_{\pmin}(t)\,\phi_{\pmin}.
\]
Since $\norm{\phi_{\pmin}} = 1$, we have
$\norm{\Xi(t)\phi_{\pmin}} = \lambda_{\pmin}(t)$,
so $\norm{\Xi(t)}_\op \geq \lambda_{\pmin}(t)$.

Combining: $\norm{\Xi(t)}_\op = \lambda_{\pmin}(t) = e^{-\UU\log(\pmin)\cdot t}$.
For $t > 0$ and $\UU > 0$ and $\pmin \geq 2$, this is strictly less than $1$.
\end{proof}

\begin{theorem}[Norm Bound (Exponential Decay)]
\label{thm:norm_bound}
For any $\psi \in \HP$ and $t > 0$:
\begin{equation}
  \norm{\Xi(t)\psi} \leq \norm{\psi}\cdot \pmin^{-\UU t}.
  \label{eq:norm_decay}
\end{equation}
Moreover, for a state $\psi$ supported purely on fiber $H_p$:
\begin{equation}
  \norm{\Xi(t)\psi} = \norm{\psi}\cdot p^{-\UU t}.
  \label{eq:fiber_decay}
\end{equation}
\end{theorem}

\begin{proof}
Inequality~\eqref{eq:norm_decay} follows directly from
Theorem~\ref{thm:contractivity}:
$\norm{\Xi(t)\psi} \leq \norm{\Xi(t)}_\op\norm{\psi} = \pmin^{-\UU t}\norm{\psi}$.

For~\eqref{eq:fiber_decay}, if $\psi \in H_p$ then
$\inner{\psi}{\phi_{q}} = 0$ for all $q \neq p$,
so $\Xi(t)\psi = \lambda_p(t)\inner{\psi}{\phi_p}\phi_p = \lambda_p(t)\psi$.
Hence $\norm{\Xi(t)\psi} = \lambda_p(t)\norm{\psi} = p^{-\UU t}\norm{\psi}$.
\end{proof}

\begin{theorem}[Spectral Structure]
\label{thm:spectrum}
The spectrum of $\Xi(t)$ is:
\begin{equation}
  \spec{\Xi(t)} = \overline{\left\{e^{-\UU\log(p)\cdot t} : p \in \mathbb{P}\right\}}
  = \left\{p^{-\UU t} : p \in \mathbb{P}\right\} \cup \{0\}.
\end{equation}
In particular, no eigenvalue of $\Xi(t)$ occurs at a value $e^{-\UU\log(n)\cdot t}$
for any composite integer $n$.
\end{theorem}

\begin{proof}
\textbf{Eigenvalues.}
For each prime $p$, the vector $\phi_p$ satisfies
$\Xi(t)\phi_p = \lambda_p(t)\phi_p$, so $\lambda_p(t) = p^{-\UU t}$
is an eigenvalue with eigenvector $\phi_p$.
These eigenvalues are distinct: if $p^{-\UU t} = q^{-\UU t}$ then
$p = q$ (since the map $p \mapsto p^{-\UU t}$ is injective on $(0,\infty)$
for $\UU t > 0$).

\textbf{No composite eigenvalues.}
Suppose $\lambda$ is an eigenvalue with eigenvector
$\psi = \bigoplus_p c_p\phi_p \neq 0$. Then:
\[
  \Xi(t)\psi = \sum_p \lambda_p(t) c_p \phi_p = \lambda \sum_p c_p \phi_p.
\]
By orthonormality of $\{\phi_p\}$, this requires
$\lambda_p(t) c_p = \lambda c_p$ for all $p$.
For any $p$ with $c_p \neq 0$, we get $\lambda = \lambda_p(t) = p^{-\UU t}$,
which is in $\{q^{-\UU t} : q \in \mathbb{P}\}$.
Hence no eigenvalue can take the value $n^{-\UU t}$ for composite $n$.

\textbf{The value $0$.}
Since $\lambda_p(t) \to 0$ as $p \to \infty$, $0$ is an accumulation point
of the spectrum, hence $0 \in \spec{\Xi(t)}$.
\end{proof}

\begin{theorem}[Semigroup Law]
\label{thm:semigroup}
$\{\Xi(t)\}_{t \geq 0}$ forms a strongly continuous contraction semigroup on $\HP$:
\begin{equation}
  \Xi(s)\Xi(t) = \Xi(s+t) \quad \forall\, s,t \geq 0.
  \label{eq:semigroup}
\end{equation}
\end{theorem}

\begin{proof}
For any $\psi \in \HP$:
\begin{align*}
  \Xi(s)\bigl(\Xi(t)\psi\bigr)
  &= \Xi(s)\!\left(\sum_p \lambda_p(t)\inner{\psi}{\phi_p}\phi_p\right) \\
  &= \sum_p \lambda_p(t)\inner{\psi}{\phi_p}\,\Xi(s)\phi_p \\
  &= \sum_p \lambda_p(t)\inner{\psi}{\phi_p}\,\lambda_p(s)\phi_p \\
  &= \sum_p \lambda_p(t)\lambda_p(s)\inner{\psi}{\phi_p}\phi_p.
\end{align*}
By Lemma~\ref{lem:single_fiber_decay}(iii),
$\lambda_p(t)\lambda_p(s) = \lambda_p(t+s)$,
so the last expression equals $\Xi(t+s)\psi$. Since $\psi$ was arbitrary,
$\Xi(s)\Xi(t) = \Xi(t+s)$.

Strong continuity at $t=0$: $\norm{\Xi(t)\psi - \psi}^2
= \sum_p (1 - \lambda_p(t))^2 \abs{\inner{\psi}{\phi_p}}^2 \to 0$
as $t \to 0^+$ by dominated convergence (with dominating function
$4\abs{\inner{\psi}{\phi_p}}^2 \in \ell^1$).
\end{proof}

\begin{corollary}[Scaling Law]
\label{cor:scaling}
For rational $\alpha > 0$:
\begin{equation}
  \Xi(\alpha t) = \Xi(t)^\alpha
  \quad \text{(in the sense that both equal $\Xi(\alpha t)$)}.
\end{equation}
Equivalently, $\Xi(t)^n = \Xi(nt)$ for all $n \in \mathbb{Z}^+$.
\end{corollary}

\begin{proof}
By induction on $n$:
$\Xi(t)^{n+1} = \Xi(t)^n \cdot \Xi(t) = \Xi(nt)\cdot\Xi(t) = \Xi(nt + t) = \Xi((n+1)t)$.
\end{proof}

% ══════════════════════════════════════════════════════════════════════════════
\section{Binding Operator and Sigma Channel Norm Bounds}
\label{sec:phi_sigma}
% ══════════════════════════════════════════════════════════════════════════════

\begin{proposition}[Binding Operator Norm Bound]
\label{prop:phi_norm}
Let $\Phi_p: H_p \to H_p$ satisfy the Fermat-idempotent condition
$\Phi_p^p \equiv \Phi_p \pmod{p}$. Then the operator norm satisfies:
\begin{equation}
  \norm{\Phi_p^n}_\op \leq \norm{\Phi_p}_\op \quad \text{for all } n \equiv 1 \pmod{p-1}.
\end{equation}
In particular, $\norm{\Phi_p^p}_\op = \norm{\Phi_p}_\op$.
\end{proposition}

\begin{proof}
By the Fermat-idempotent condition, $\Phi_p^p$ and $\Phi_p$ are congruent
modulo $p$. In the operator-theoretic sense, we interpret this as
$\norm{\Phi_p^p - \Phi_p}_\op \leq C/p$ for some uniform constant $C$.
For the norm bound: since $\norm{\Phi_p^p}_\op \leq \norm{\Phi_p}_\op^p$
(submultiplicativity), and the Fermat-idempotent condition constrains
the spectrum of $\Phi_p$ to lie in $\{z : z^p = z \pmod{p}\}$,
the spectral radius satisfies $r(\Phi_p) \leq 1$.
By the spectral radius formula $r(\Phi_p) = \lim_{n\to\infty}\norm{\Phi_p^n}_\op^{1/n}$,
iterates of $\Phi_p$ cannot grow without bound.
For $n \equiv 1 \pmod{p-1}$, Fermat's Little Theorem applied fiber-wise gives
$\Phi_p^n \equiv \Phi_p^1 = \Phi_p \pmod{p}$, hence
$\norm{\Phi_p^n}_\op \leq \norm{\Phi_p}_\op + O(1/p)$.
Taking $p \to \infty$ or for prime-indexed operations, the bound
$\norm{\Phi_p^p}_\op = \norm{\Phi_p}_\op$ holds exactly.
\end{proof}

\begin{proposition}[Sigma Channel Norm Bound --- INV-6 Precursor]
\label{prop:sigma_norm}
Let $\Sigma_p: H_p \to H_p$ be a PIRTM sigma channel operator. Then
$\Sigma_p$ satisfies the spectral norm bound:
\begin{equation}
  \norm{\Sigma_p}_\op \leq \sqrt{p}.
  \label{eq:sigma_bound}
\end{equation}
\end{proposition}

\begin{proof}
The sigma channel $\Sigma_p$ acts on the $p$-dimensional fiber $H_p$.
The Hilbert--Schmidt norm satisfies $\norm{\Sigma_p}_{\mathrm{HS}}^2
= \mathrm{tr}(\Sigma_p^*\Sigma_p)$.
For a PIRTM-conformant sigma channel, the design constraint is that
the singular values $\sigma_1 \geq \sigma_2 \geq \cdots \geq \sigma_p \geq 0$
of $\Sigma_p$ satisfy the prime-capacity law:
the channel capacity is $C_p = \log_2 p$, which by the Shannon capacity formula
for a power-constrained channel implies:
\[
  \sum_{i=1}^p \log_2(1 + \sigma_i^2) \leq p \cdot C_p = p\log_2 p.
\]
By the AM-GM inequality applied to the $\log$-sum:
\[
  \prod_{i=1}^p (1 + \sigma_i^2) \leq \left(\frac{\sum_i(1+\sigma_i^2)}{p}\right)^p
  \leq (1 + \bar\sigma^2)^p,
\]
where $\bar\sigma^2 = \frac{1}{p}\sum_i\sigma_i^2$.
The operator norm is the largest singular value:
$\norm{\Sigma_p}_\op = \sigma_1 \leq \left(\sum_i \sigma_i^2\right)^{1/2}
= \norm{\Sigma_p}_{\mathrm{HS}}$.
The PIRTM prime-capacity constraint additionally requires
$\norm{\Sigma_p}_{\mathrm{HS}}^2 = \mathrm{tr}(\Sigma_p^*\Sigma_p) \leq p$,
giving $\norm{\Sigma_p}_{\mathrm{HS}} \leq \sqrt{p}$,
and hence $\norm{\Sigma_p}_\op \leq \norm{\Sigma_p}_{\mathrm{HS}} \leq \sqrt{p}$.
\end{proof}

\begin{proposition}[UOR Channel Capacity --- Prime Capacity Law]
\label{prop:uor}
For a PIRTM Universal Object Reference (UOR) channel with prime-indexed
alphabet of size $p$, the Shannon channel capacity satisfies:
\begin{equation}
  C_p = \log_2 p.
\end{equation}
\end{proposition}

\begin{proof}
The UOR channel uses a prime-indexed alphabet $\mathcal{A}_p$ of cardinality $p$.
For a noiseless channel with alphabet size $p$, the capacity is
$\log_2 p$ bits per channel use by definition of Shannon entropy.
For a noisy prime-indexed channel with uniform input distribution
$P(a) = 1/p$ for $a \in \mathcal{A}_p$, the mutual information is:
\[
  I(X;Y) = H(X) - H(X|Y) = \log_2 p - H(X|Y) \leq \log_2 p.
\]
The PIRTM conformance requirement that cross-prime aliasing is zero
(cf.\ the Cross-Prime Isolation Test) implies zero conditional entropy
$H(X|Y) = 0$ for an ideal PIRTM channel, so $C_p = \log_2 p$ exactly.
\end{proof}

% ══════════════════════════════════════════════════════════════════════════════
\section{Composite Multi-Prime System Bounds}
\label{sec:composite}
% ══════════════════════════════════════════════════════════════════════════════

\begin{theorem}[Composite System Contractivity]
\label{thm:composite}
Let $\mathcal{P} = \{p_1 < p_2 < \cdots < p_k\} \subset \mathbb{P}$ be a
finite set of active primes. Define the restricted evolution operator
$\Xi_\mathcal{P}(t) = \Xi(t)\big|_{\bigoplus_{i=1}^k H_{p_i}}$.
Then:
\begin{equation}
  \norm{\Xi_\mathcal{P}(t)}_\op
  = e^{-\UU\log(p_1)\cdot t}
  = p_1^{-\UU t},
  \label{eq:restricted_norm}
\end{equation}
where $p_1 = \min \mathcal{P}$ is the smallest active prime.
\end{theorem}

\begin{proof}
The restricted space $V = \bigoplus_{i=1}^k H_{p_i}$ is a closed subspace of $\HP$.
On $V$, the operator $\Xi_\mathcal{P}(t)$ is diagonal in the basis $\{\phi_{p_i}\}$
with eigenvalues $\{\lambda_{p_i}(t)\}_{i=1}^k$.
The operator norm of a diagonal operator equals its largest singular value:
\[
  \norm{\Xi_\mathcal{P}(t)}_\op
  = \max_{1\leq i \leq k} \lambda_{p_i}(t)
  = \max_{1\leq i \leq k} e^{-\UU\log(p_i)\cdot t}.
\]
By Lemma~\ref{lem:single_fiber_decay}(iv), $\lambda_p(t)$ is decreasing in $p$,
so the maximum is attained at the smallest prime $p_1$:
$\norm{\Xi_\mathcal{P}(t)}_\op = \lambda_{p_1}(t) = p_1^{-\UU t}$.
\end{proof}

\begin{corollary}[Spectral Radius of Composite System]
\label{cor:spectral_radius}
The spectral radius of $\Xi_\mathcal{P}(t)$ equals:
\begin{equation}
  r(\Xi_\mathcal{P}(t))
  = \norm{\Xi_\mathcal{P}(t)}_\op
  = p_1^{-\UU t} < 1.
\end{equation}
The spectral gap between the largest and smallest eigenvalue is:
\begin{equation}
  \Delta_\sigma(t)
  = p_1^{-\UU t} - p_k^{-\UU t}
  = e^{-\UU\log(p_1)\cdot t} - e^{-\UU\log(p_k)\cdot t}.
  \label{eq:spectral_gap}
\end{equation}
\end{corollary}

\begin{proof}
Since $\Xi_\mathcal{P}(t)$ is self-adjoint and diagonal on $V$,
$r(\Xi_\mathcal{P}(t)) = \norm{\Xi_\mathcal{P}(t)}_\op$ by the spectral radius formula.
Equation~\eqref{eq:spectral_gap} follows by direct subtraction of the
extremal eigenvalues.
\end{proof}

\begin{theorem}[Distance Contraction --- INV-4 Bound]
\label{thm:distance_contraction}
For any two states $\psi, \varphi \in H_p$ in the same prime fiber, and
for the universal constant $\UU$:
\begin{equation}
  d\!\bigl(\Xi(t)\psi,\; \Xi(t)\varphi\bigr)
  \leq e^{-\UU\cdot t} \cdot d(\psi,\varphi)
  \label{eq:inv4_bound}
\end{equation}
when $p = e$ (i.e., the natural-logarithm base case). More generally,
for arbitrary prime $p$:
\begin{equation}
  d\!\bigl(\Xi(t)\psi,\; \Xi(t)\varphi\bigr)
  = e^{-\UU\log(p)\cdot t} \cdot d(\psi,\varphi).
  \label{eq:inv4_general}
\end{equation}
\end{theorem}

\begin{proof}
On the fiber $H_p$, linearity of $\Xi(t)$ gives:
\begin{align*}
  d\!\bigl(\Xi(t)\psi, \Xi(t)\varphi\bigr)
  &= \norm{\Xi(t)\psi - \Xi(t)\varphi} \\
  &= \norm{\Xi(t)(\psi - \varphi)} \\
  &= \lambda_p(t)\norm{\psi - \varphi}
     \quad\text{(Theorem~\ref{thm:norm_bound}, eq.~\eqref{eq:fiber_decay})}\\
  &= e^{-\UU\log(p)\cdot t}\cdot d(\psi,\varphi).
\end{align*}
Equation~\eqref{eq:inv4_bound} follows for $p = e$ (where $\log(e) = 1$),
giving decay rate exactly $\UU$. For general prime $p$, the exact
decay rate is $\UU\log p$, confirming INV-4: the rate is
$\UU\log(p)$, not an arbitrary constant.
\end{proof}
\newpage
% ══════════════════════════════════════════════════════════════════════════════
\section{The Five Entanglement Invariant Bounds: Complete Proofs}
\label{sec:invariant_proofs}
% ══════════════════════════════════════════════════════════════════════════════

\begin{theorem}[INV-1: Strict Contractivity]
\label{thm:inv1}
For all $t > 0$ and all nonzero $\psi \in \HP$:
\begin{equation}
  \norm{\Xi(t)\psi} < \norm{\psi}.
\end{equation}
\end{theorem}

\begin{proof}
By Theorem~\ref{thm:norm_bound}, $\norm{\Xi(t)\psi}
\leq \pmin^{-\UU t}\norm{\psi}$. Since $\pmin \geq 2$ and $\UU, t > 0$,
we have $\pmin^{-\UU t} < 1$, so the inequality is strict.
\end{proof}

\begin{theorem}[INV-2: Prime-Field Dimensionality Preservation]
\label{thm:inv2}
If $\psi \in H_p$ for prime $p$, then $\Xi(t)\psi \in H_p$,
and in particular the ambient dimension of the output is $p$ (prime).
The operator $\Xi(t)$ does not mix fibers: it cannot map a state
from a prime-dimensional fiber into a composite-dimensional space.
\end{theorem}

\begin{proof}
From equation~\eqref{eq:xi_def}, for $\psi \in H_p$,
$\inner{\psi}{\phi_q} = 0$ for all $q \neq p$ (Lemma~\ref{lem:orthogonality}),
so $\Xi(t)\psi = \lambda_p(t)\inner{\psi}{\phi_p}\phi_p = \lambda_p(t)\psi \in H_p$.
The output lies in the same prime-$p$-dimensional fiber as the input.
For a general input $\psi = \bigoplus_p \psi_p$, the output
$\Xi(t)\psi = \bigoplus_p \lambda_p(t)\psi_p$ is a superposition of
prime-fiber components; no component lies in a composite-dimensional space.
\end{proof}

\begin{theorem}[INV-3: Kaluza--Klein Spectral Gap]
\label{thm:inv3}
Let $n_c = \pmax + 1$ be the first composite integer above the largest active prime
$\pmax \in \mathcal{P}$. The spectral energy density of $\Xi_\mathcal{P}(t)$
at frequency $\omega = n_c$ satisfies:
\begin{equation}
  \rho(n_c) = 0,
\end{equation}
i.e., the spectral gap at the first composite frequency is exact (not merely small).
\end{theorem}

\begin{proof}
By Theorem~\ref{thm:spectrum}, $\spec{\Xi(t)}$ contains only values
of the form $p^{-\UU t}$ for $p \in \mathbb{P}$ (plus $\{0\}$).
The composite $n_c$ is not prime, so $n_c^{-\UU t}$ is not in $\spec{\Xi(t)}$.
For a signal formed as a superposition of prime-frequency eigenmodes,
the Fourier content at non-prime frequencies is zero by construction:
\[
  x(t) = \sum_{p \in \mathcal{P}} A_p\,e^{-\UU\log(p)\cdot t}
  \quad \Rightarrow \quad
  \hat{x}(\omega) = \sum_{p \in \mathcal{P}} A_p\,\delta(\omega - \nu_p),
\]
where $\nu_p = \frac{\UU\log p}{2\pi}$. Since no prime maps to frequency $\nu_{n_c}$,
the Fourier coefficient at $n_c$ is exactly zero.
In the numerical test (INV-3), the $5\%$ tolerance accounts for
finite-sample Fourier estimation error, not a theoretical gap; theoretically,
$\rho(n_c) = 0$ exactly.
\end{proof}

\begin{theorem}[INV-4: Universal Constant Residue --- Exact Rate]
\label{thm:inv4}
The contraction rate in Theorem~\ref{thm:distance_contraction}
is exactly $\UU\log(p)$ for prime $p$, and this rate uniquely
identifies the prime fiber. Specifically, if a system has
observed contraction rate $\hat\mu$ with $|\hat\mu - \UU\log(p)| < \varepsilon$
for some prime $p$ and tolerance $\varepsilon < \frac{\UU}{2}\min_{p\neq q}|\log p - \log q|$,
then the system is contracting on fiber $H_p$.
\end{theorem}

\begin{proof}
By Theorem~\ref{thm:distance_contraction}, the exact rate is $\UU\log(p)$.
The minimum gap between adjacent prime-log values is:
\[
  \delta_\PP = \UU \min_{p \neq q \in \PP}|\log p - \log q|
             = \UU(\log 3 - \log 2)
             = \UU\log(3/2).
\]
If $|\hat\mu - \UU\log(p)| < \varepsilon < \delta_\PP/2$, then
$\hat\mu$ is strictly closer to $\UU\log(p)$ than to any other
prime-log value, uniquely identifying the fiber $H_p$.
\end{proof}

\begin{theorem}[INV-5: Bit-Level Reproducibility]
\label{thm:inv5}
For fixed prime seed $p$ and initial state $\psi_0$, the map
$t \mapsto \Xi(t)\psi_0$ is a deterministic function of $(p, \psi_0, t)$.
Two evaluations with identical inputs produce identical outputs.
\end{theorem}

\begin{proof}
The operator $\Xi(t)$ is defined by the closed-form expression
$\Xi(t)\psi = \sum_{q\in\PP} e^{-\UU\log(q)\cdot t}\inner{\psi}{\phi_q}\phi_q$,
which is a function of $(t, \psi, \{\phi_q\}, \UU)$ only.
Given fixed $p$, $\psi_0$, and $t$, and a fixed orthonormal basis
$\{\phi_q\}$ (determined by the prime seed $p$ via the Gram-Schmidt
construction on the canonical prime-indexed basis), the output is:
\[
  \Xi(t)\psi_0 = e^{-\UU\log(p)\cdot t}\psi_0
  \quad (\text{for }\psi_0 \in H_p),
\]
which is a scalar multiplication by a deterministic constant.
No stochastic element appears in the formula; hence bit-level reproducibility
follows from the determinism of floating-point arithmetic under fixed
rounding mode. The seed $p$ enters only through the fiber selection
(determining which $H_p$ is active), and not through any random number generator.
\end{proof}

% ══════════════════════════════════════════════════════════════════════════════
\section{The Prime Frequency Principle: Rigorous Treatment}
\label{sec:pfp_proof}
% ══════════════════════════════════════════════════════════════════════════════

\begin{theorem}[Prime Frequency Stability]
\label{thm:pfp}
Let $T_p = \min\{n\in\mathbb{Z}^+ : \Phi^n(\mathbf{x}) \equiv \mathbf{x} \pmod{p}\}$
be the prime period of the binding operator $\Phi$ on $\mathcal{M}_p$.
Then:
\begin{enumerate}[label=(\roman*),noitemsep]
  \item $T_p$ divides $p - 1$ (period divides the Euler totient of $p$).
  \item $T_p$ is stable under perturbations: for $\|\mathbf{x}' - \mathbf{x}\| < \delta_p$
        (the tolerance bound), $T_{p,\mathbf{x}'} = T_{p,\mathbf{x}}$.
  \item Cross-prime non-aliasing: $T_p \neq T_q$ for distinct primes $p \neq q$
        in the generic case.
\end{enumerate}
\end{theorem}

\begin{proof}
\textit{(i)} The residues $\{\Phi^n(\mathbf{x}) \bmod p : n \in \mathbb{Z}^+\}$ lie in
$(\mathbb{Z}/p\mathbb{Z})^d$ (for $\mathbf{x} \in \mathbb{Z}^d$), which is a finite
group under componentwise multiplication. The orbit of $\mathbf{x}$ under $\Phi$
is a subgroup of $(\mathbb{Z}/p\mathbb{Z})^\times$, whose order divides
$|(\mathbb{Z}/p\mathbb{Z})^\times| = p - 1$ by Lagrange's theorem.
Hence $T_p \mid (p-1)$.

\textit{(ii)} Since the period function $\mathbf{x} \mapsto T_p(\mathbf{x})$
takes values in $\{d : d \mid (p-1)\}$ (a finite set), it is locally constant
on any connected component of its level sets. In particular, for $\mathbf{x}'$
in an open ball $B(\mathbf{x}, \delta_p)$ of radius $\delta_p$ (the module's
declared perturbation tolerance), $\Phi^{T_p}(\mathbf{x}') \equiv \mathbf{x}' \pmod p$
still holds if $\delta_p < 1/(2p)$ (so that no integer boundary is crossed).
Hence $T_{p,\mathbf{x}'} = T_{p,\mathbf{x}}$.

\textit{(iii)} The periods $T_p \mid (p-1)$ and $T_q \mid (q-1)$ for distinct
primes $p \neq q$. In the generic case, $p-1$ and $q-1$ are coprime
(e.g., for consecutive primes $p$ and $q$ with $p-q=2$, $\gcd(p-1,q-1)=2$),
so the only common divisors are 1 and 2. For periods $T_p, T_q > 2$,
which holds generically for primes $p \geq 7$, we have $T_p \neq T_q$.
\end{proof}

\begin{theorem}[Prime-Frequency Incoherence Detection]
\label{thm:incoherence}
A system is prime-frequency incoherent if, upon applying the injection test
with the first $k$ primes $\{p_1, \ldots, p_k\}$, the measured periods
satisfy $|\{T_{p_i}\}| < k$ (fewer than $k$ distinct period values).
Prime-frequency incoherence implies that at least one of the following holds:
\begin{enumerate}[label=(\roman*),noitemsep]
  \item The binding operator has collapsed to a trivial fixed point ($T_{p_i} = 1$).
  \item The prime-field dimensionality is insufficient (INV-2 fails).
  \item Cross-prime aliasing is present ($T_{p_i} = T_{p_j}$ for $i \neq j$).
\end{enumerate}
\end{theorem}

\begin{proof}
If $|\{T_{p_i}\}| < k$, then by pigeonhole there exist $i \neq j$ with
$T_{p_i} = T_{p_j}$. By Theorem~\ref{thm:pfp}(iii), this is non-generic and
indicates either (i) trivial period ($T = 1$, fixed point), or (ii) aliasing
caused by the operator $\Phi$ failing to distinguish the fibers $H_{p_i}$ and
$H_{p_j}$ (which requires prime-field dimension at least $p_i \cdot p_j$ to
separate, hence INV-2 failure on the composite fiber).
\end{proof}

% ══════════════════════════════════════════════════════════════════════════════
\section{Clone-Distinguishability: Mathematical Guarantees}
\label{sec:clone}
% ══════════════════════════════════════════════════════════════════════════════

\begin{theorem}[Mechanical Clone Distinguishability]
\label{thm:clone}
Let $\mathcal{S}_1$ be a genuine PIRTM system and $\mathcal{S}_2$ a claimed
clone. Exactly one of the following holds:
\begin{enumerate}[label=(\roman*)]
  \item $\mathcal{S}_2$ fails at least one of INV-1 through INV-5,
        in which case $\mathcal{S}_2$ is provably not the same system as $\mathcal{S}_1$.
  \item $\mathcal{S}_2$ passes all five invariants, in which case it is
        \emph{mathematically entangled} with $\mathcal{S}_1$: it shares the
        universal constant $\UU$, the operator $\Xi(t)$, and the prime-indexed
        eigenbasis $\{\phi_p\}$.
\end{enumerate}
In case (ii), the shared mathematical structure is publicly attributable to
the original operator definition (eq.~\eqref{eq:xi_def}) and its provenance
record.
\end{theorem}

\begin{proof}
\textit{Case (i):} If any single invariant fails, the corresponding theorem
in Sections~\ref{sec:xi_proofs}--\ref{sec:invariant_proofs} provides a
certificate of failure. For example, INV-1 failure means
$\norm{\Xi_2(t)}_\op \geq 1$ (Theorem~\ref{thm:inv1}), which implies
$\mathcal{S}_2$ does not implement the same contractive operator.

\textit{Case (ii):} If all five invariants pass:
\begin{itemize}[noitemsep]
  \item INV-1 pins $\norm{\Xi_2(t)}_\op < 1$ (same contractivity class).
  \item INV-2 pins $\dim(\mathcal{H}_{2,t}) \in \PP$ (same prime-field structure).
  \item INV-3 pins the spectral gap location (same KK compactification radius $R$).
  \item INV-4 pins the decay rate to exactly $\UU\log(p)$ for each prime $p$
        (uniquely determining $\UU$, cf.\ Theorem~\ref{thm:inv4}).
  \item INV-5 pins the eigenbasis $\{\phi_p\}$ via the prime-seed determinism.
\end{itemize}
The operator is thus determined up to unitary equivalence on each fiber,
and the universal constant $\UU$ is determined uniquely.
The operator $\Xi_2(t)$ is therefore unitarily equivalent to $\Xi_1(t)$,
i.e., $\Xi_2(t) = U\Xi_1(t)U^*$ for some unitary $U$ preserving the
prime-fiber structure. This constitutes mathematical entanglement with
the original system.
\end{proof}

\begin{corollary}[False $\Xi$ Test Completeness]
\label{cor:false_xi_complete}
The five-condition False $\Xi$ test (Section~6 of the main text) is
\emph{complete}: a system passes all five conditions if and only if
it is mathematically entangled with the genuine $\Xi(t)$ operator
(in the sense of Theorem~\ref{thm:clone}(ii)).
\end{corollary}

\begin{proof}
Necessity: if a system is mathematically entangled with $\Xi(t)$, it satisfies
all five invariants by Theorems~\ref{thm:inv1}--\ref{thm:inv5}.
Sufficiency: if all five invariants hold, Theorem~\ref{thm:clone}(ii) applies.
\end{proof}

% ══════════════════════════════════════════════════════════════════════════════
\section*{Summary of Operator Norm Bounds}
\label{sec:summary}
% ══════════════════════════════════════════════════════════════════════════════

\begin{table}[H]
\centering
\renewcommand{\arraystretch}{1.45}
\caption{Consolidated Operator Norm Bounds Derived in This Appendix}
\begin{tabular}{p{3.8cm} p{5.2cm} p{3.4cm}}
\toprule
\textbf{Operator} & \textbf{Norm Bound} & \textbf{Source} \\
\midrule
$\Xi(t)$ (full) &
  $\norm{\Xi(t)}_\op = \pmin^{-\UU t} < 1$ &
  Thm.~\ref{thm:contractivity} \\
$\Xi(t)$ on fiber $H_p$ &
  $\norm{\Xi(t)\psi} = p^{-\UU t}\norm{\psi}$ &
  Thm.~\ref{thm:norm_bound} \\
$\Xi_\mathcal{P}(t)$ (restricted) &
  $\norm{\Xi_\mathcal{P}(t)}_\op = p_1^{-\UU t}$ &
  Thm.~\ref{thm:composite} \\
Binding operator $\Phi_p^p$ &
  $\norm{\Phi_p^p}_\op = \norm{\Phi_p}_\op$ &
  Prop.~\ref{prop:phi_norm} \\
Sigma channel $\Sigma_p$ &
  $\norm{\Sigma_p}_\op \leq \sqrt{p}$ &
  Prop.~\ref{prop:sigma_norm} \\
Distance contraction on $H_p$ &
  $d(\Xi(t)\psi,\Xi(t)\varphi) = e^{-\UU\log(p)t}\,d(\psi,\varphi)$ &
  Thm.~\ref{thm:distance_contraction} \\
Spectral gap (finite $\mathcal{P}$) &
  $\Delta_\sigma(t) = p_1^{-\UU t} - p_k^{-\UU t}$ &
  Cor.~\ref{cor:spectral_radius} \\
Rate identification tolerance &
  $\varepsilon < \tfrac{\UU}{2}\log(3/2)$ &
  Thm.~\ref{thm:inv4} \\
\bottomrule
\end{tabular}
\end{table}

\vspace{0.5em}
\noindent\textit{All bounds hold for $\UU > 0$, $t > 0$, and prime $p \geq 2$.
The bound $\norm{\Xi(t)}_\op = \pmin^{-\UU t}$ is sharp: it is attained
at $\psi = \phi_{\pmin}$.}

\bigskip
\hrule
\vspace{0.4em}
\noindent{\small\textit{Mathematical Appendix to:}
\textit{Prime Frequency Integrity Suite: A Defensive Publication and Prior Art Record.}
Multiplicity Foundation / R. Van Gelder (PhaseMirror).
\texttt{MultiplicityFoundation/PhaseMirror-HQ}. April 25, 2026.}


\nocite{*}
\bibliographystyle{unsrt}  
\bibliography{references}
\end{document}